\documentclass[11pt]{article}
\usepackage{fontspec}
\setmainfont{Times New Roman}
\setsansfont{Arial}
\setmonofont{Consolas}
\usepackage{amsmath,amssymb,mathtools}
\usepackage{geometry}
\usepackage{microtype}
\geometry{a4paper,margin=2.35cm}
\setlength{\parindent}{1.25em}
\setlength{\parskip}{0pt}
\makeatletter
\renewcommand\footnotesize{\@setfontsize\footnotesize{7.7}{8.7}\selectfont}
\makeatother
\emergencystretch=100em
\allowdisplaybreaks
\sloppy
\hbadness=10000
\vbadness=10000
\hfuzz=2pt
\vfuzz=10pt
\raggedbottom

\begin{document}

% Унит: Папер31 Сецтион06 но. 4 опенинг в001. Дирецт Интерславиц Латин транслатион фром the Герман финал-аудитед соурце слице, соурце линес 358--362 / сегмент P31-С0080, експандед агаинст принтед сцан/ОЦР витнесс паге 98. Фоотноте цоунтер цонтинуес афтер Папер31 Сецтион06 но. 1.
\setcounter{footnote}{27}

4. \emph{Представјенье следа, нормы и дискриминанта чрез конјуговане елементы в случају полној разложимости првого рода.} Нехај по \S~3, 3 при полној разложимости првого рода
\[
        \mathfrak R[Q]=Qe_1+\cdots+Qe_n
\]
јест представјенье како директна сума колц ранга један, то јест поль изоморфных к $Q$, и нехај $Q$ јест најменше разширјено полье, в ктором тако представјенье јест можно. Нехај $n$ компонентов од $\mathfrak R$ сут дане чрез $K_i e_i$, где всако $K_i$ јест подполье од $Q$ и $\mathfrak R$ јест хомоморфно к $K_i$; треба изразити след, норму и дискриминант (взходно к главному класу) чрез елементы из $K_i$.

\end{document}

